\pagestyle{empty}

\chapter*{Resumen}
\thispagestyle{empty}

% Entrada en el índice sin número de página
\addtocontents{toc}{
    \protect\contentsline{chapter}{Resumen}{}{}
}

Se desarrolló un modelo computacional de un acelerador
lineal de radiofrecuencia con tubos de deriva, de estructura
tipo Alvarez, que opera a una frecuencia de
$352.2\,\mathrm{MHz}$ y permite estudiar la aceleración de
protones desde $3$ hasta $50\,\mathrm{MeV}$. El diseño se
estableció a partir de la condición de sincronismo
$L_n=\beta_n\lambda$ y de la ganancia energética por cavidad,
$\Delta E=qV\cos\phi$, con $qV=1.0\,\mathrm{MeV}$.
Para una fase síncrona de $\phi_s=-30^\circ$, se obtuvieron
55 celdas y una longitud total de $10.59\,\mathrm{m}$,
correspondientes a un gradiente medio de aceleración de
$4.44\,\mathrm{MeV/m}$. Posteriormente, se simuló un paquete
de 3000 protones con una desviación estándar inicial de fase
de $15.0^\circ$ y una dispersión relativa de energía de
\SI{0.5}{\percent}. Los resultados muestran que, para
$\phi_s=-30^\circ$, el \SI{99.3}{\percent} del paquete
permanece agrupado longitudinalmente, mientras que la
desviación estándar de fase disminuye de $15.0^\circ$
a $5.3^\circ$, en concordancia con el amortiguamiento
adiabático. En contraste, para la operación en la cresta
de radiofrecuencia ($\phi_s=0^\circ$), únicamente el
\SI{24.1}{\percent} de las partículas alcanza la energía
final de diseño. Asimismo, la comparación entre la
simulación y la aproximación lineal para las oscilaciones
de sincrotrón presenta un avance de fase de $25.7^\circ$
por celda, frente a los $26.4^\circ$ predichos
teóricamente. En cuanto a la sincronización, se encontró
que un error del \SI{2}{\percent} en la frecuencia de
radiofrecuencia permite conservar el \SI{94.5}{\percent}
de las partículas capturadas, mientras que un error del
\SI{4}{\percent} reduce esta fracción al
\SI{55.5}{\percent}. Finalmente, se observa que los
electrones con velocidades cercanas a la de la luz
($\beta\approx 1$) no presentan oscilaciones de fase
apreciables.

\vspace{6pt}

\noindent
\hspace*{0.5in}
\textit{Palabras clave:}
LINAC, acelerador lineal, estructura de Alvarez, estabilidad de fase,
dinámica longitudinal de haces, oscilaciones de sincrotrón,
simulación computacional.


% ------------------------------------------------------------
% ABSTRACT
% ------------------------------------------------------------

\clearpage

\begin{otherlanguage}{english}

\chapter*{Abstract}
\thispagestyle{empty}

% Entry in the table of contents without page number
\addtocontents{toc}{
    \protect\contentsline{chapter}{Abstract}{}{}
}

A computational model of a radio-frequency drift-tube
linear accelerator (DTL) with an Alvarez-type structure,
operating at a frequency of $352.2\,\mathrm{MHz}$, was
developed to study proton acceleration from $3$ to
$50\,\mathrm{MeV}$. The design was based on the synchronism
condition $L_n=\beta_n\lambda$ and the energy gain per
accelerating gap, $\Delta E=qV\cos\phi$, with
$qV=1.0\,\mathrm{MeV}$. For a synchronous phase of
$\phi_s=-30^\circ$, the resulting design required 55 cells
and a total length of $10.59\,\mathrm{m}$, corresponding to
an average accelerating gradient of $4.44\,\mathrm{MeV/m}$.
Subsequently, a bunch of 3,000 protons was simulated with
an initial phase standard deviation of $15.0^\circ$ and a
relative energy spread of \SI{0.5}{\percent}. The results
show that, for $\phi_s=-30^\circ$, \SI{99.3}{\percent}
of the particle bunch remains longitudinally captured,
while the phase standard deviation decreases from
$15.0^\circ$ to $5.3^\circ$, consistent with adiabatic
damping. In contrast, when operating at the
radio-frequency crest ($\phi_s=0^\circ$), only
\SI{24.1}{\percent} of the particles reach the final
design energy. Furthermore, comparison of the simulated
synchrotron phase oscillations with the linear approximation
yields a phase advance of $25.7^\circ$ per cell, compared
with the theoretical prediction of $26.4^\circ$. Regarding
synchronization, a \SI{2}{\percent} radio-frequency error
allows \SI{94.5}{\percent} of the particles to remain
captured, whereas a \SI{4}{\percent} error reduces this
fraction to \SI{55.5}{\percent}. Finally, electrons
traveling at velocities close to the speed of light
($\beta\approx 1$) exhibit no appreciable phase oscillations.

\vspace{6pt}

\noindent
\hspace*{0.5in}
\textit{Keywords:}
linear accelerator, Alvarez structure, phase stability,
longitudinal beam dynamics, synchrotron oscillations,
computational simulation.

\end{otherlanguage}
